\section{The next three routes}
\label{sec:decision-tree}

The audit reduces the next block to two positive inputs and one
negative-transfer alternative:
\begin{equation}
\boxed{
\begin{gathered}
\text{Route I: determinant-energy survival}\\
\text{Route II: distinguished-zero-mode cancellation}\\
\text{Route III: exact nonidentifiability or no-go theorem.}
\end{gathered}}
\label{eq:three-routes}
\end{equation}
Routes I and II can proceed in parallel, but their exponents meet only
through \eqref{eq:zero-slack-audit}.  Route III becomes appropriate
when a branch-specific obstruction is proved or both positive routes
are shown to require the same unresolved arithmetic input.

\subsection{Route I: the literal dictionary and energy survival}

The target is
\begin{equation}
 D_X\ge X^{-o(1)}Q_X^3J_X^2,
\label{eq:routeI-natural-energy}
\end{equation}
or, after an independent zero-mode gain, the more general inequality
\(\lambda_D\le2\eta_Z\).  There are two nonexclusive subroutes.

\paragraph{Route I-A: completion, reassembly, and dictionary.}
This certificate route has six ordered tasks.

\begin{enumerate}[label=(I\arabic*),leftmargin=2.7em]
\item Generate a complete literal manifest at every audited scale,
retaining the fixed \(h_0\), all native keys, actual support, literal
phases, and the global raw normalization.
\item Construct \(\cJ_X\) from the exact determinant label and archive
every multiplier and provenance contribution in each fiber.
\item Run the TPC-77 component census on the actual nonzero incidence
graph.  Use TPC-78 forest, one-defect, or motif formulas only on
components whose literal graph satisfies the corresponding
hypotheses.
\item Use TPC-79 only when the anchor and direct variables are
separable syntheses into the same target with the physical costs.
Otherwise delete the hybrid branch and keep the direct route.
\item Use TPC-80 only on an actual scalar row graph.  Audit both
sectors, and record exact switching equality rather than interpreting
it as cancellation.
\item Prove a positive pre-bin survival factor, a lower modulus
\(s_X\), and residual separation strong enough that
\eqref{eq:dictionary-energy-audit} yields the desired
\(\lambda_D\).
\end{enumerate}

\paragraph{Route I-B: direct determinant-energy analysis.}
Starting from the already defined literal coefficient
\(a_X=(A_{C_X}(n))_n\), expand or estimate
\[
 D_X=\sum_n|A_{C_X}(n)|^2
\]
directly.  This route does not logically require a completion
manifest, a pre-bin carrier, or a lower-modulus argument.  It must
still preserve the fixed \(h_0\), the complete definition of
\(A_{C_X}(n)\), the actual post-aggregation support, and the global
normalization.  The exact expansion should separate manifestly
nonnegative diagonal mass from signed equal-determinant
cross-correlations.  If a natural-scale lower bound is equivalent to
an unresolved fixed-shift Chowla-like correlation, that equivalence is
a valid Route-III negative-transfer theorem rather than a proof of
energy survival.

\begin{proposition}[Route-I branch rules]
\label{prop:routeI-rules}
The following conclusions are local to their branch.
\begin{enumerate}[label=(\alph*),leftmargin=2.2em]
\item Different targets or coupled variables stop the TPC-79 formula,
but do not disprove direct synthesis.
\item A nonforest component stops the forest shortcut, but leaves the
general component SOCP and dual program available.
\item A tree, or exact switching equality in a dominant scalar
component, stops phase-frustration saving on that component, but does
not decide amplitude completion or determinant cancellation.
\item A dual lower bound exceeding the allowed cost disproves only
the declared completion class and normalization.
\item If the currently available Route-I-A certificate yields only
\(\lambda_D^{\rm cert}>0\) while \(\eta_Z=0\), that certificate cannot
close the present endpoint transfer.  A sharper certificate or a
direct Route-I-B estimate remains available.  Only a proof that the
literal carrier necessarily incurs positive loss is a branch
obstruction.
\end{enumerate}
\end{proposition}

\begin{proof}
Items (a)--(d) follow from the hypotheses and targets of
\cref{prop:tpc79-admissibility,cor:tree-two-sector,thm:no-teleportation}.
Item (e) follows from \cref{cor:natural-energy-audit} together with
the branch-locality principle: failure of one sufficient certificate
does not prove failure of a sharper lower bound.
\end{proof}

\subsection{Route II: the distinguished zero mode}

The target is an estimate
\begin{equation}
 |Z_X|\le X^{-\eta_Z+o(1)}Q_X^2,
\qquad
 2\eta_Z\ge\lambda_D,
\label{eq:routeII-target}
\end{equation}
with strict inequality preferred because it supplies reserve.  A
valid input must concern:
\begin{enumerate}[label=(Z\arabic*),leftmargin=2.7em]
\item the same prescribed \(h_0\ne0\);
\item the actual matched coefficient \(A_{C_X}(n)\);
\item the distinguished frequency \(r=0\), not merely generic
\(r\ne0\);
\item the complete physical block and the original normalization; and
\item the original uniformity, with no unreturned averaging over
shift, frequency, slope, modulus, or carrier family.
\end{enumerate}

\begin{proposition}[Results that do not close Route II]
\label{prop:not-routeII}
None of the following, without a proved localization or transfer back
to the literal packet, establishes \eqref{eq:routeII-target}:
\begin{enumerate}[label=(\roman*),leftmargin=2.2em]
\item an all-but-exceptional-frequency estimate;
\item an average over \(h\), slopes, moduli, or frozen families;
\item squarefree-support density or an absolute-value Perron bound;
\item cancellation on one fiber followed by absolute summation in an
outer variable; or
\item a nonzero-frequency dispersion theorem.
\end{enumerate}
\end{proposition}

\begin{proof}
Each item changes one of (Z1)--(Z5) or discards the sign interaction
defining \(Z_X\).  It therefore does not bound the statistic in
\eqref{eq:routeII-target}.
\end{proof}

Expanding \(D_X\) or \(Z_X\) may expose a fixed-shift two-affine
M\"obius correlation.  If the desired bound is shown to be equivalent
to a currently unproved Chowla-type estimate, that estimate must be
listed as a new arithmetic hypothesis.  Averaged Chowla results do
not automatically give the required literal form
\citep{MatomakiRadziwillTao2015,Tao2016LogChowla}.

\subsection{Route III: a negative transfer theorem}

If Routes I and II both reduce to the same unresolved signed
correlation, continuing to rename the input creates no progress.
An honest alternative is a paper of the form
\begin{quote}
\emph{Why Completion Geometry Does Not Control the Distinguished
Frequency: Nonidentifiability and Stop Theorems for a Fixed-Shift
Transfer Architecture.}
\end{quote}
Its rigorous content can include:
\begin{enumerate}[label=(N\arabic*),leftmargin=2.7em]
\item identical support, absolute coefficients, and energy with
different zero modes;
\item complete loss of raw energy under determinant binning;
\item failure of anchor--direct composition without a common target;
\item absence of scalar phase saving on trees or saturated sectors;
\item failure of scalar Perron strictness to imply determinant
flatness; and
\item exact branch-specific dual obstructions for independent,
provenance, hybrid, or holonomy models.
\end{enumerate}
Such a result is an \(\Lzero\) no-go theorem.  It becomes an
\(\Lone\) negative diagnostic only when its countermodel or dual
witness is attached to the literal fixed-\(h_0\) interface.  It is
not a proof of the general sieve parity barrier and does not disprove
a twin-prime strategy outside the declared architecture.

\subsection{Operational decision tree}

\begin{equation}
\begin{aligned}
&\text{choose a determinant-energy branch:}\\
&\quad\begin{cases}
\text{I-A:} & \text{build the literal manifest and \(\cJ_X\), then
estimate \(s_X,\delta_{\rm reasm},\Gamma_{\rm pre},r_X\);}\\
\text{I-B:} & \text{expand the literal \(D_X\) directly and isolate
its signed cross-correlations.}
\end{cases}\\[0.4em]
&\text{a physical determinant-energy exponent \(\lambda_D\) known?}\\
&\quad\begin{cases}
\text{no:} & \text{continue Route I or identify the signed
correlation obstruction;}\\
\text{yes:}& \text{attack the same coefficient's \(Z_X\) in Route II.}
\end{cases}\\[0.4em]
&\text{compare the two exponents:}\\
&\quad\begin{cases}
2\eta_Z>\lambda_D:& \text{positive reserve; audit the full physical
loss ledger;}\\
2\eta_Z=\lambda_D:& \text{flatness threshold only; no extra
loss inside this flatness transfer;}\\
2\eta_Z<\lambda_D:& \text{stop the current TPC-32 endpoint route.}
\end{cases}
\end{aligned}
\label{eq:operational-tree}
\end{equation}

Even the first case in the last comparison reaches only the
conditional matched-shell theorem.  A full-block conclusion still
requires the exact cover and remainder control isolated in TPC-18
\citep{WangTPC18}.
